\section{Definitions and Measurement}
\label{sec:framework}

\paragraph{Load-bearing components.} A multi-agent organization is defined
by its agents, interaction order, message-passing graph, context-visibility
policy and role contracts. Prohibitive (\textsc{must not}) clauses, such as
instructing the Planner not to state a final answer, are what establish a
division of labor. Because these clauses are only natural-language prompts,
the realized organization need not match the specification. We call a
component \textbf{load-bearing} if removing or altering it reproducibly
degrades task performance under matched call and information budgets, and we apply it separately to first answers and to repairs
(Appendix~\ref{app:loadbearing}).

\paragraph{Functional collapse.} The first deliberation cycle
\emph{collapses} if all three agents independently assert their own final
answers. The pipeline then operates as three sequential solvers rather than
as a division of labor. This is our main measure of functional
specialization. It does not apply to the structurally enforced condition,
whose schema has no answer field. Appendix~\ref{app:measures} defines the
secondary measures.

\paragraph{First pass and recovery.} A \emph{trajectory} is a team's run on
one problem. \emph{First-pass accuracy} grades the first answer the team
submits to the evaluator. It need not come in the first cycle: a
role-structured team may first deliberate internally, without external
feedback. A trajectory that never
submits is graded once on its final answer when the budget runs out, and
that grade is its first pass. Trajectories whose first graded answer is
wrong are \emph{first-pass wrong} (failed, for short), and \emph{recovery} is the share of them
that finish correct. Because recovery conditions on a post-treatment
variable (whether the first answer failed, which the intervention can
change), recovery contrasts are descriptive~\cite{montgomery2018conditioning}.
Final accuracy is reported over all problems.
